%------------------------------------------------------------------------------
% Security Analysis of ESAFL with KGC Zero-Ciphertext Compensation
% Compatible with IEEE single-column format
%------------------------------------------------------------------------------
\section{Security Analysis}
\label{sec:security}

In this section, we present a formal security analysis of the proposed \textsc{ECPFL} scheme under four adversarial models: semi-honest clients, client collusion, a malicious server, and client-server collusion. All proofs rely on the standard Decisional Ring Learning With Errors (D-RLWE) hardness assumption and the collision-resistance of the homomorphic hash function.

%------------------------------------------------------------------------------
\subsection{Threat Model and Assumptions}
\label{ssec:threat-model}

\noindent\textbf{System Participants.}
The protocol involves three types of entities:
\begin{itemize}
    \item \textbf{Key Generation Center (KGC)}: Trusted but minimally involved. It generates MK-CKKS key pairs, derives the public-polynomial sequence $\{a^{(t)}\}_{t=1}^{T}$, pre-computes all compensation ciphertexts, and then goes permanently offline. The KGC never observes any training data or gradients.
    \item \textbf{Server $\mathcal{S}$}: \emph{Malicious}. It may deviate from the protocol, attempt to infer individual client gradients, tamper with aggregation results, forge virtual clients, or collude with other participants.
    \item \textbf{Clients $\mathcal{C}_i$ ($i \in [N]$)}: Each client may be semi-honest (honest-but-curious) or malicious. Semi-honest clients follow the protocol exactly but try to infer information about other clients from protocol transcripts.
\end{itemize}

\noindent\textbf{Cryptographic Assumptions.}
\begin{assumption}[RLWE Hardness]
\label{ass:rlwe}
For security parameter $\lambda$, the Decisional-RLWE problem over $R_q = \mathbb{Z}_q[x]/(x^n+1)$ is hard. That is, for $a \gets R_q$, $s \gets \chi$, $e \gets \chi$, the tuple $(a, a \cdot s + e)$ is computationally indistinguishable from a uniformly random tuple $(a, u) \gets R_q^2$.
\end{assumption}

\begin{assumption}[Collision-Resistant Homomorphic Hash]
\label{ass:hash}
The hash function $\mathcal{H}: \{0,1\}^* \to \mathbb{G}$ satisfies collision-resistance and pre-image resistance, where $\mathbb{G}$ is a cyclic group of prime order $p'$.
\end{assumption}

\begin{assumption}[Trusted KGC]
\label{ass:kgc}
The KGC honestly executes key generation and compensation ciphertext construction, securely erases all private key material after going offline, and is never compromised by any adversary.
\end{assumption}

%------------------------------------------------------------------------------
\subsection{Security Against Semi-Honest Clients}
\label{ssec:semi-honest}

Under the semi-honest client threat model, we require that no client $i$ can infer any other client $j \neq i$'s local gradient from its protocol view, even though client $i$ ultimately learns the global aggregated gradient.

Let $f$ denote the ideal functionality that takes all local gradients $\mathcal{G} = \{\mathbf{g}_1, \dots, \mathbf{g}_N\}$ and outputs the aggregated gradient $\mathbf{g}_{\text{agg}}^{(t)} = \sum_{i \in \mathcal{S}^{(t)}} \mathbf{g}_i^{(t)}$. Let $\mathsf{View}_i(\lambda, \mathcal{G})$ denote the view of client $i$ during protocol execution, including all received messages (global model $\mathbf{w}^{(t)}$, hash anchors $H_{\text{pre}}^{(t)}$, and the decrypted cumulative gradient $\mathbf{\Delta}_{\text{pre}}^{(i)}$), as well as its own random coins and internal state.

\begin{theorem}[IND-CPA Security against Semi-Honest Clients]
\label{thm:semi-honest}
Under Assumption~\ref{ass:rlwe}, for any proper subset of semi-honest clients, there exists a PPT simulator $\mathsf{Sim}_i$ such that for all PPT adversaries $\mathcal{A}$:
\begin{equation}
    \bigl| \Pr[\mathcal{A}(\mathsf{View}_i(\lambda, \mathcal{G})) = 1] - \Pr[\mathcal{A}(\mathsf{Sim}_i(\mathbf{g}_i, f(\mathcal{G}))) = 1] \bigr| \leq \mathsf{negl}(\lambda).
\end{equation}
\end{theorem}

\begin{proof}
The view of client $i$ contains only two sources of information about other clients: (1) the ciphertexts uploaded by other clients to the server, which client $i$ never receives directly; and (2) the decrypted aggregated gradient $\mathbf{\Delta}_{\text{pre}}^{(i)}$.

For (1), by the IND-CPA security of RLWE (Assumption~\ref{ass:rlwe}), for any two plaintexts $\mathbf{m}_0, \mathbf{m}_1$:
\begin{equation}
    c_{j,\mathrm{up}}^{(t)} = \bigl[ a^{(t)} \cdot s_j + \mathsf{CRTStack}(\mathbf{g}_j) + e_j \bigr]_{R_q}
\end{equation}
is computationally indistinguishable from a uniformly random element of $R_q$. Even if client $i$ were to obtain $c_{j,\mathrm{up}}^{(t)}$ (e.g., via server collusion analyzed in \S\ref{ssec:collusion}), decryption would still require $s_j$, which remains unknown.

For (2), the simulator $\mathsf{Sim}_i$ simply outputs the ideal functionality result $f(\mathcal{G}) = \mathbf{\Delta}_{\text{pre}}^{(t)}$, which is exactly what client $i$ is supposed to learn. Since the protocol reveals nothing beyond this aggregation, the simulated view is perfectly indistinguishable from the real view.
\end{proof}

\noindent\textbf{Key Observation.} The aggregated gradient is the \emph{intended output} of federated learning---every client needs it to update its local model. Therefore, the security lower bound is: client $i$ cannot learn anything \emph{beyond} the aggregation result.

%------------------------------------------------------------------------------
\subsection{Resistance to Client Collusion}
\label{ssec:client-collusion}

Consider a collusion set $\mathcal{T} \subset [N]$ with $|\mathcal{T}| = t < N$. The colluding clients may share their local gradients $\{\mathbf{g}_i\}_{i \in \mathcal{T}}$, private key fragments $\{s_i\}_{i \in \mathcal{T}}$, and all observed protocol transcripts. Their goal is to infer the individual gradient of any non-colluding client $k \notin \mathcal{T}$.

\begin{theorem}[Collusion Lower Bound]
\label{thm:collusion}
In an additive aggregation scheme, if $|\mathcal{T}| < |\mathcal{S}^{(t)}| - 1$, the colluding set $\mathcal{T}$ cannot recover the exact gradient of any non-colluding client $k \notin \mathcal{T}$ with non-negligible probability.
\end{theorem}

\begin{proof}
At round $t$, the decrypted aggregated gradient is:
\begin{equation}
    \mathbf{\Delta}_{\mathrm{pre}}^{(t)} = \sum_{i \in \mathcal{S}^{(t)}} \mathbf{g}_i^{(t)} + \tilde{E}^{(t)},
\end{equation}
where $\tilde{E}^{(t)}$ is the residual decryption noise, which is eliminated during exact integer reconstruction via CRT.

The colluding clients know their own gradient sum:
\begin{equation}
    \mathbf{G}_{\mathcal{T}} = \sum_{i \in \mathcal{T} \cap \mathcal{S}^{(t)}} \mathbf{g}_i^{(t)}.
\end{equation}

By subtracting their known contribution, they obtain:
\begin{equation}
    \mathbf{G}_{\mathcal{S} \setminus \mathcal{T}} = \mathbf{\Delta}_{\mathrm{pre}}^{(t)} - \mathbf{G}_{\mathcal{T}} = \sum_{i \in \mathcal{S}^{(t)} \setminus \mathcal{T}} \mathbf{g}_i^{(t)}.
\end{equation}

This is a linear system with $|\mathcal{S}^{(t)}| - |\mathcal{T}|$ unknown vectors in $\mathbb{R}^d$ and $d$ equations. When $|\mathcal{S}^{(t)}| - |\mathcal{T}| \geq 2$, the system is \emph{underdetermined}, admitting infinitely many solutions. Therefore, the colluding clients cannot uniquely determine any individual $\mathbf{g}_k$ for $k \notin \mathcal{T}$.
\end{proof}

\noindent\textbf{Pairwise Blinding Mask Enhancement.}
If pairwise blinding masks $\mathbf{b}_{i,j}$ are enabled (satisfying $\sum_i w_i \mathbf{b}_{i,j} = \mathbf{0}$), clients encrypt $\tilde{\mathbf{g}}_i = \mathbf{g}_i + \mathbf{b}_i$ where $\mathbf{b}_i = \sum_{j \neq i} \mathbf{b}_{i,j}$. The colluding set $\mathcal{T}$ only knows $\{\mathbf{b}_{i,j}\}_{i,j \in \mathcal{T}}$ and cannot eliminate mask terms involving $j \notin \mathcal{T}$. Thus, even when $|\mathcal{S}^{(t)}| - |\mathcal{T}| = 1$, the mask protects the sole non-colluding client.

%------------------------------------------------------------------------------
\subsection{Security Against a Malicious Server}
\label{ssec:malicious-server}

A malicious server $\mathcal{S}^*$ may attempt: (1) decryption attacks to recover individual gradients; (2) tampering attacks to manipulate the global model; (3) replay attacks using old ciphertexts; (4) compensation ciphertext substitution; and (5) Sybil attacks by forging virtual clients.

\subsubsection{Resistance to Decryption Attacks}

\begin{theorem}[Server UnDecryptability]
\label{thm:server-decrypt}
Under Assumptions~\ref{ass:rlwe} and~\ref{ass:kgc}, a malicious server cannot recover any individual client's plaintext gradient from the aggregated ciphertext $C_j^{(t)}$ with non-negligible probability.
\end{theorem}

\begin{proof}
The aggregated ciphertext is:
\begin{equation}
    C_j^{(t)} = \Bigl[ a^{(t)} \cdot S + \sum_{i \in \mathcal{S}^{(t)}} m_{i,j}^{(t)} + E^{(t)} \Bigr]_{R_q},
\end{equation}
where $S = \sum_{\ell=1}^{N} s_\ell$ is the global aggregated secret key. The server possesses $C_j^{(t)}$ and the public polynomial $a^{(t)}$, but \emph{does not possess any} $s_\ell$.

To remove the key-mask term $a^{(t)} \cdot S$, the server would need to compute:
\begin{equation}
    \tilde{m} = \bigl[ C_j^{(t)} - a^{(t)} \cdot S \bigr]_{R_q}.
\end{equation}

Since $S$ is unknown, this is impossible. Furthermore, by the RLWE pseudorandomness (Assumption~\ref{ass:rlwe}), for any $S \gets \chi$, the tuple $(a^{(t)}, a^{(t)} \cdot S + E^{(t)})$ is computationally indistinguishable from uniform over $R_q^2$. Therefore:
\begin{equation}
    \Pr_{a,S,E}\bigl[\mathcal{A}(a, aS+E) = S\bigr] \leq \mathsf{negl}(\lambda),
\end{equation}
for any PPT adversary $\mathcal{A}$.
\end{proof}

\subsubsection{Resistance to Tampering Attacks}

If the server tampers with the aggregated ciphertext as $\tilde{C}_j^{(t)} = C_j^{(t)} + \Delta$ for some malicious perturbation $\Delta$, the client verifies integrity before decryption:
\begin{equation}
    \mathsf{Verify}\bigl(\mathcal{H}(\tilde{C}_j^{(t)}), H_{\mathrm{pre},j}^{(t)}\bigr) \stackrel{?}{=} 1.
\end{equation}

By the collision-resistance of $\mathcal{H}$ (Assumption~\ref{ass:hash}):
\begin{equation}
    \Pr\bigl[\mathcal{H}(\tilde{C}_j^{(t)}) = \mathcal{H}(C_j^{(t)}) \mid \tilde{C}_j^{(t)} \neq C_j^{(t)}\bigr] \leq \mathsf{negl}(\lambda).
\end{equation}
Thus, tampering is detected with overwhelming probability.

\subsubsection{Resistance to Compensation Ciphertext Substitution}

The server may attempt to replace a KGC-generated compensation ciphertext with a self-generated one. The compensation structure is:
\begin{equation}
    c_{\mathrm{comp},j}^{(t,k)} = \bigl[ a^{(t)} \cdot s_k + e_{\mathrm{comp},j}^{(t,k)} \bigr]_{R_q}.
\end{equation}

Since the server does not know $s_k$, any forgery $\tilde{c} = [a^{(t)} \cdot \tilde{s} + \tilde{e}]_{R_q}$ introduces an irremovable residual term $a^{(t)}(\tilde{s} - s_k)$ during decryption:
\begin{equation}
    \tilde{m} = \bigl[ m_{\mathrm{true}} + a^{(t)}(\tilde{s} - s_k) + \tilde{E} \bigr]_{R_p},
\end{equation}
which destroys the integer boundary required for CRT exact reconstruction. The client detects this as a decoding failure.

\subsubsection{Resistance to Sybil Attacks}

The server cannot forge new client identities because each legitimate client requires a private key fragment $s_i$ distributed by the KGC. The server does not participate in $s_i$ generation, and the KGC is offline and trusted (Assumption~\ref{ass:kgc}).

%------------------------------------------------------------------------------
\subsection{Client-Server Collusion Boundary}
\label{ssec:collusion}

We now consider the most dangerous scenario where a client subset $\mathcal{T} \subsetneq [N]$ colludes with the malicious server $\mathcal{S}^*$.

\begin{theorem}[Partial Collusion Security]
\label{thm:partial-collusion}
Even if the server colludes with a proper subset $\mathcal{T} \subsetneq [N]$ of clients, as long as $|\mathcal{T}| < N$, the colluding parties cannot recover the plaintext gradient of any non-colluding client $k \notin \mathcal{T}$.
\end{theorem}

\begin{proof}
The colluding parties possess $\{s_i\}_{i \in \mathcal{T}}$ and all ciphertexts from the server-side view. Decryption requires the global secret:
\begin{equation}
    S = \sum_{i \in \mathcal{T}} s_i + \sum_{k \notin \mathcal{T}} s_k = S_{\mathcal{T}} + S_{\neg\mathcal{T}}.
\end{equation}

Let $S_{\mathcal{T}} = \sum_{i \in \mathcal{T}} s_i$ (known to the colluders) and $S_{\neg\mathcal{T}} = \sum_{k \notin \mathcal{T}} s_k$ (unknown). The colluders can compute:
\begin{equation}
    C_j^{(t)} - a^{(t)} S_{\mathcal{T}} = \bigl[ a^{(t)} S_{\neg\mathcal{T}} + \sum_{i \in \mathcal{S}} m_{i,j}^{(t)} + E^{(t)} \bigr]_{R_q}.
\end{equation}

The remaining term $a^{(t)} S_{\neg\mathcal{T}}$ is protected by RLWE (Assumption~\ref{ass:rlwe}). Since $S_{\neg\mathcal{T}}$ is the sum of at least one independent uniform secret (as $|\mathcal{T}| < N$), $a^{(t)} S_{\neg\mathcal{T}} + E$ remains pseudorandom. Therefore, the colluding parties cannot distinguish the aggregation from random noise.
\end{proof}

\subsubsection{Worst-Case Analysis: Single Online Client}

Consider the extreme case where $|\mathcal{S}^{(t)}| = 1$ and the sole online client $i$ colludes with the server. The aggregated ciphertext degenerates to:
\begin{equation}
    C_j^{(t)} = c_{i,j,\mathrm{up}}^{(t)} \oplus \bigoplus_{k \neq i} c_{\mathrm{comp},j}^{(t,k)} = \Bigl[ a^{(t)} S + m_{i,j}^{(t)} + E_{\mathrm{total}} \Bigr]_{R_q}.
\end{equation}

The colluding parties know $s_i$ and can compute:
\begin{equation}
    C_j^{(t)} - a^{(t)} s_i = \Bigl[ a^{(t)} \underbrace{\sum_{k \neq i} s_k}_{S_{\neg i}} + m_{i,j}^{(t)} + E_{\mathrm{total}} \Bigr]_{R_q}.
\end{equation}

However, $S_{\neg i} = \sum_{k \neq i} s_k$ remains an RLWE secret. The colluding parties cannot distinguish $a^{(t)} S_{\neg i} + m_{i,j}^{(t)} + E$ from $a^{(t)} S_{\neg i} + E'$ (another noise sample). Thus, \emph{even in this extreme case, a single client colluding with the server still cannot decrypt}.

\subsubsection{Total Collusion: Security Collapse}

When $\mathcal{T} = [N]$ (all clients collude with the server), $S_{\neg\mathcal{T}} = 0$, and the colluding parties fully know $S$. Then:
\begin{equation}
    \tilde{m} = \bigl[ C_j^{(t)} - a^{(t)} S \bigr]_{R_q} = \sum_{i \in \mathcal{S}} m_{i,j}^{(t)} + E^{(t)}.
\end{equation}
Security collapses completely---the colluding parties recover the exact aggregation. This is the \emph{theoretical lower bound} of secure aggregation: \textbf{no scheme can resist collusion by all participants}.

%------------------------------------------------------------------------------
\subsection{Comparison with Standard Secure Aggregation}
\label{ssec:comparison}

Table~\ref{tab:security-comparison} compares the collusion resistance of \textsc{ECPFL} with the classical Bonawitz et al. secure aggregation (SecAgg).

\begin{table}[htbp]
\centering
\caption{Collusion Resistance Comparison}
\label{tab:security-comparison}
\begin{tabular}{lcc}
\toprule
\textbf{Collusion Scenario} & \textbf{SecAgg} & \textbf{ESAFL + KGC} \\
\midrule
Server alone & Cannot decrypt & Cannot decrypt (no $S$) \\
$t < N/2$ clients only & Secure & Secure (underdetermined system) \\
Client + Server & Secure if $t < N/2$ & Secure unless $t = N$ \\
All $N$ clients + Server & Collapse & Collapse (theoretical lower bound) \\
\bottomrule
\end{tabular}
\end{table}

The key advantage of ESAFL + KGC compensation is that \textbf{server collusion with any proper subset of clients remains insufficient for decryption}. The adversary must collect \emph{all $N$} private key fragments, whereas in SecAgg, the server colluding with a simple majority ($> N/2$) suffices.

%------------------------------------------------------------------------------
\subsection{Security of Compensation Ciphertexts}
\label{ssec:compensation-security}

\subsubsection{Zero-Plaintext Indistinguishability}

The KGC-generated compensation ciphertexts have zero plaintext but fresh noise:
\begin{equation}
    c_{\mathrm{comp},j}^{(t,k)} = \bigl[ a^{(t)} s_k + e_{\mathrm{comp},j}^{(t,k)} \bigr]_{R_q}.
\end{equation}

\begin{lemma}[Compensation Indistinguishability]
\label{lem:comp-ind}
For any two distinct dropped clients $k_1, k_2$, their compensation ciphertexts $c_{\mathrm{comp},j}^{(t,k_1)}$ and $c_{\mathrm{comp},j}^{(t,k_2)}$ are computationally indistinguishable, assuming $s_{k_1}, s_{k_2}$ are drawn independently from $\chi$.
\end{lemma}

\begin{proof}
By Assumption~\ref{ass:rlwe}, $(a^{(t)}, a^{(t)} s_k + e)$ is pseudorandom for any $s_k \gets \chi$. Since $s_{k_1}$ and $s_{k_2}$ are independent and identically distributed, the two ciphertext tuples follow identical distributions.
\end{proof}

Therefore, compensation ciphertexts leak no information about the identity of dropped clients.

\subsubsection{Server Cannot Distinguish Real from Compensation}

Comparing an active client's ciphertext $c_{i,j,\mathrm{up}}^{(t)} = [a^{(t)} s_i + m_{i,j}^{(t)} + e_i]_{R_q}$ with a compensation ciphertext $c_{\mathrm{comp},j}^{(t,k)} = [a^{(t)} s_k + e_k]_{R_q}$:
\begin{itemize}
    \item From the server's perspective, both follow the RLWE distribution.
    \item Even if the server knows which clients dropped out (and thus which ciphertexts are compensations), it cannot extract any information from the zero-plaintext compensations.
    \item The server \emph{cannot distinguish} whether a given ciphertext is a real update or a compensation without the private key.
\end{itemize}

%------------------------------------------------------------------------------
\subsection{Summary}
\label{ssec:summary}

The security of ESAFL with KGC zero-ciphertext compensation rests on four pillars:
\begin{enumerate}
    \item \textbf{RLWE Hardness}: Guarantees ciphertext indistinguishability for any party without the global secret.
    \item \textbf{Distributed Multi-Key Structure}: The global key $S = \sum_{i=1}^{N} s_i$ is split across $N$ clients; no single point of compromise exists.
    \item \textbf{Minimal Trusted KGC}: The KGC only participates in initialization, never touches training data, and goes offline permanently after provisioning.
    \item \textbf{Unforgeable Compensation Structure}: The server cannot forge compensations without knowing the missing clients' $s_k$.
\end{enumerate}

Under semi-honest clients, client collusion, malicious servers, and partial client-server collusion, the scheme provides provable security. Security collapses only under the theoretical lower bound where \emph{all $N$ clients and the server fully collude}---a scenario no aggregation scheme can resist.
